%%
%% Datei: paper74_lonely_runner_de.tex
%% Autor: Michael Fuhrmann
%% Beschreibung: Die Lonely-Runner-Vermutung (Deutsche Version)
%% Erstellt: 2026-03-12
%% lastModified: 2026-03-12
%% Build: 164
%%
\documentclass[12pt,a4paper]{amsart}

\usepackage{amsmath,amssymb,amsthm,mathtools,enumitem,booktabs,array,hyperref}
\usepackage[utf8]{inputenc}
\usepackage[T1]{fontenc}
\usepackage[ngerman]{babel}
\usepackage{geometry}
\geometry{margin=2.5cm}

%% Theorem-Umgebungen (Englisch, Basis)
\newtheorem{theorem}{Theorem}[section]
\newtheorem{lemma}[theorem]{Lemma}
\newtheorem{corollary}[theorem]{Corollary}
\newtheorem{proposition}[theorem]{Proposition}
\theoremstyle{definition}
\newtheorem{definition}[theorem]{Definition}
\newtheorem{example}[theorem]{Example}
\theoremstyle{remark}
\newtheorem{remark}[theorem]{Remark}
\newtheorem{conjecture}[theorem]{Conjecture}

%% Deutsche Theorem-Umgebungen
\newtheorem{satz}[theorem]{Satz}
\newtheorem{vermutung}[theorem]{Vermutung}
\newtheorem{korollar}[theorem]{Korollar}
\newtheorem{lemma_de}[theorem]{Lemma}
\newtheorem{bemerkung}[theorem]{Bemerkung}
\newtheorem{definition_de}[theorem]{Definition}

%% Kürzel
\newcommand{\Z}{\mathbb{Z}}
\newcommand{\R}{\mathbb{R}}
\newcommand{\Q}{\mathbb{Q}}
\newcommand{\N}{\mathbb{N}}
\newcommand{\T}{\mathbb{T}}
\newcommand{\norm}[1]{\left\|#1\right\|}

%% Abstract VOR \maketitle
\begin{abstract}
Die Lonely-Runner-Vermutung, ursprünglich von Wills \cite{Wills1967} vorgeschlagen
und von Cusick und Pomerance \cite{CusickPomerance1984} neu formuliert, fragt,
ob auf einer Kreisbahn des Umfangs 1 mit $n$ Läufern konstanter ganzzahliger
Geschwindigkeit jeder Läufer zu irgendeinem Zeitpunkt einen Abstand von mindestens
$1/(n+1)$ zu allen anderen Läufern hat ("einsam wird"). Wir geben eine
Gesamtdarstellung der Geschichte, der äquivalenten Formulierungen (diophantische
Approximation, Sichtverstopfungsproblem), der Beweise für $n \leq 7$ Läufer, der
Verbindungen zur chromatischen Zahl von Kneser-Graphen sowie der wichtigsten
offenen Probleme. Die Vermutung ist für $n \leq 7$ Läufer bewiesen
(Barajas--Serra \cite{BarajasSerra2008}), im Allgemeinen jedoch offen.
\end{abstract}

\title{Die Lonely-Runner-Vermutung}

\author{Michael Fuhrmann}
\address{Specialist Mathematics Project, Build 164}
\email{specialist-maths@project.local}
\date{2026-03-12}

\maketitle

\tableofcontents

%% ============================================================
\section{Einleitung}
%% ============================================================

Stellen wir uns $n+1$ Läufer auf einer Kreisbahn des Umfangs 1 vor. Läufer $i$
(für $i = 0, 1, \ldots, n$) läuft mit konstanter ganzzahliger Geschwindigkeit $v_i$,
wobei $0 = v_0 < v_1 < v_2 < \cdots < v_n$. Wir sagen, Läufer 0 ist zum
Zeitpunkt $t$ \emph{einsam} (lonely), wenn sein Abstand zu jedem anderen Läufer
mindestens $1/(n+1)$ beträgt:
\[
  \|v_i t\| := \min_{k \in \Z} |v_i t - k| \geq \frac{1}{n+1}
  \quad \text{für alle } i = 1, \ldots, n,
\]
wobei $\|\cdot\|$ den Abstand zur nächsten ganzen Zahl bezeichnet.

\begin{vermutung}[Lonely-Runner-Vermutung \cite{Wills1967,CusickPomerance1984}]
\label{verm:lrc}
Für alle $n \geq 1$ und alle positiven ganzen Zahlen $v_1 < v_2 < \cdots < v_n$
gibt es eine reelle Zahl $t$ mit
\[
  \|v_i t\| \geq \frac{1}{n+1} \quad \text{für alle } i = 1, \ldots, n.
\]
\end{vermutung}

Die Schranke $1/(n+1)$ ist bestmöglich: Für $v_i = i$ ($i = 1, \ldots, n$) kann
der ruhende Läufer einen Abstand von höchstens $1/(n+1)$ zu jedem anderen Läufer
haben, und Gleichheit wird bei $t = 1/(n+1)$ erreicht.

\subsection{Historische Anmerkungen}

Das Problem wurde 1967 von Wills \cite{Wills1967} in einer anderen Formulierung
(Sichtverstopfung; vgl.\ Abschnitt~\ref{sec:sicht}) eingeführt, 1984 von Cusick
und Pomerance \cite{CusickPomerance1984} neu formuliert und von Bienia et al.\
\cite{BieniaEtal1998} als "Lonely Runner" bezeichnet.

%% ============================================================
\section{Äquivalente Formulierungen}\label{sec:aequiv}
%% ============================================================

\subsection{Formulierung als diophantische Approximation}

\begin{satz}[Äquivalenz]\label{satz:dioph}
Die Lonely-Runner-Vermutung ist äquivalent dazu: Für beliebige $n$ von Null
verschiedene ganze Zahlen $w_1, \ldots, w_n$ gibt es $t \in \R$ mit
$\|w_i t\| \geq 1/(n+1)$ für alle $i$.
\end{satz}

\begin{proof}
Durch Subtraktion der Geschwindigkeit von Läufer 0 von allen anderen wird die
Einsamkeitsbedingung zu $\|(v_i - v_0)t\| \geq 1/(n+1)$ für alle $i \geq 1$.
Die Substitution $w_i = v_i - v_0$ liefert die äquivalente Formulierung.
\end{proof}

\subsection{Sichtverstopfungsformulierung}\label{sec:sicht}

\begin{definition_de}[Sichtverstopfungsproblem]
Platziere einen Einheitswürfel $[-1/2, 1/2]^n$ an jedem Gitterpunkt von $\Z^n$.
Eine Gerade durch den Ursprung mit Richtungsvektor $(v_1, \ldots, v_n)$ heißt
\emph{verstopft}, wenn sie einen Würfel trifft. Die LRC ist äquivalent zur Aussage:
Jede Gerade durch den Ursprung (in Richtung $(1, v_1, \ldots, v_n)$) trifft einen
Würfel im "inneren" Sinne.
\end{definition_de}

\begin{bemerkung}
Cusick und Pomerance \cite{CusickPomerance1984} zeigten, dass die
Sichtverstopfungsformulierung für $n$-dimensionale Würfel mit Seitenlänge $1/(n+1)$
äquivalent zur LRC ist.
\end{bemerkung}

\subsection{Verbindung zur chromatischen Zahl}

\begin{satz}[LRC und zirkuläre chromatische Zahl \cite{TothCayley2001}]\label{satz:chrom}
Die Lonely-Runner-Vermutung ist äquivalent dazu: Für jedes $n$ und jede
$n$-elementige Menge $S \subset \Z \setminus \{0\}$ gilt
$\chi_c(G(S)) \geq n+1$,
wobei $G(S)$ der von $S$ erzeugte Cayley-Graph auf $\Z$ und $\chi_c$ die
zirkuläre chromatische Zahl ist.
\end{satz}

%% ============================================================
\section{Bekannte Fälle}\label{sec:bekannt}
%% ============================================================

\subsection{$n = 1$ Läufer}

\begin{satz}[$n = 1$]\label{satz:n1}
Für jede von Null verschiedene ganze Zahl $v$ gibt es $t$ mit $\|vt\| \geq 1/2$.
\end{satz}
\begin{proof}
Wähle $t = 1/(2v)$; dann ist $\|vt\| = \|1/2\| = 1/2 \geq 1/2$. \end{proof}

\subsection{$n = 2$ Läufer}

\begin{satz}[$n = 2$, Betke--Wills \cite{BetkeWills1972}]\label{satz:n2}
Für zwei positive ganze Zahlen $v_1 < v_2$ gibt es $t$ mit
$\|v_1 t\| \geq 1/3$ und $\|v_2 t\| \geq 1/3$.
\end{satz}

\begin{proof}
Die Funktion $f(t) = \min(\|v_1 t\|, \|v_2 t\|)$ muss auf $[0,1)$ ihr Maximum
annehmen. Die $2(v_1 + v_2)$ "schlechten" Intervalle (wo $f(t) < 1/3$) haben
Gesamtmaß $< 1$, sodass das gute Komplement nichtleer ist.
\end{proof}

\subsection{Allgemeiner Fall via Dirichlet}

\begin{satz}[Dirichlet-Typschranke]\label{satz:dirichlet}
Für beliebige $n$ positive ganze Zahlen $v_1, \ldots, v_n$ und jedes $\varepsilon > 0$
gibt es $t \in \R$ mit $\|v_i t\| \geq 1/(n+1) - \varepsilon$ für alle $i$.
\end{satz}

\begin{proof}
Nach Kroneckers Theorem ist die Bahn $\{(v_1 t, \ldots, v_n t) \bmod 1 : t \in \R\}$
dicht in einer abgeschlossenen Untergruppe des Torus $\T^n$. Die "schlechte" Menge
für Läufer $i$ hat Maß $2/(n+1)$. Da $n \cdot 2/(n+1) < 2$, ist das gute Komplement
nichtleer.
\end{proof}

\begin{bemerkung}
Dieser Satz gibt nur $\|v_i t\| \geq 1/(n+1) - \varepsilon$, nicht die scharfe
Schranke $1/(n+1)$. Letztere ist der Inhalt der LRC.
\end{bemerkung}

\subsection{Die Fälle $n \leq 7$}

\begin{satz}[$n \leq 7$, Barajas--Serra \cite{BarajasSerra2008}]\label{satz:n7}
Die Lonely-Runner-Vermutung ist wahr für $n \leq 7$ Läufer.
\end{satz}

Historische Progression der bewiesenen Fälle:
\begin{center}
\begin{tabular}{lll}
\toprule
\textbf{$n$ Läufer} & \textbf{Autoren} & \textbf{Jahr} \\
\midrule
$n = 1$ & Trivial & --- \\
$n = 2$ & Betke--Wills \cite{BetkeWills1972} & 1972 \\
$n = 3$ & Cusick \cite{Cusick1973} & 1974 \\
$n = 4$ & Cusick--Pomerance \cite{CusickPomerance1984} & 1984 \\
$n = 5$ & Bohman--Holzman--Kleitman \cite{BohmanHolzmanKleitman2001} & 2001 \\
$n = 6$ & Barajas--Serra \cite{BarajasSerra2008} & 2008 \\
$n = 7$ & Barajas--Serra \cite{BarajasSerra2008} & 2008 \\
\bottomrule
\end{tabular}
\end{center}

\begin{proof}[Beweisskizze für $n = 3$]
Für $n = 3$ Läufer mit Geschwindigkeiten $0 < v_1 < v_2 < v_3$ hat das Komplement
der "schlechten" Menge für jeden Läufer Maß $> 1/2$. Ein Inklusion-Exklusion-Argument
zeigt, dass der Durchschnitt der drei "guten" Mengen positives Maß hat. Für $n \leq 4$
reicht dieses Maßargument; für $n \geq 5$ sind feinere zahlentheoretische Argumente
erforderlich.
\end{proof}

%% ============================================================
\section{Arithmetische Progressionen}\label{sec:ap}
%% ============================================================

\begin{satz}[LRC für arithmetische Progressionen \cite{BieniaEtal1998}]\label{satz:lrc_ap}
Sind $v_i = i \cdot d$ für eine positive ganze Zahl $d$ und $i = 1, \ldots, n$
(d.\,h.\ die Geschwindigkeiten bilden eine arithmetische Progression), so gilt die LRC.
\end{satz}

\begin{proof}
Wähle $t = 1/(n+1)$. Dann gilt $v_i t = id/(n+1)$. Die Werte $id/(n+1)$ modulo 1
für $i = 1, \ldots, n$ liegen alle im Abstand $\geq 1/(n+1)$ von ganzen Zahlen,
wenn $\gcd(d, n+1) = 1$. Im allgemeinen Fall wähle $t$ geeignet, sodass
$td$ in einem "Einsamkeitsintervall" liegt.
\end{proof}

%% ============================================================
\section{Spektrale und Fourier-Methoden}\label{sec:spektral}
%% ============================================================

\begin{definition_de}[Spektrum einer Läuferkonfiguration]
Die "schlechte" Menge zum Zeitpunkt $t$ ist die Menge der Läufer $j$ mit
$\|v_j t\| < 1/(n+1)$. Die Vermutung behauptet, dass diese Menge für ein $t$
leer ist.
\end{definition_de}

\begin{satz}[Fourier-Kriterium]\label{satz:fourier}
Die LRC für Geschwindigkeiten $v_1, \ldots, v_n$ ist äquivalent zu:
\[
  \int_0^1 \prod_{i=1}^{n} \mathbf{1}_{[\|v_i t\| \geq 1/(n+1)]} \, dt > 0.
\]
\end{satz}

\begin{bemerkung}
Diese Fourier-Umformulierung reduziert die LRC auf eine Frage über die Positivität
einer quadratischen Form über ganzzahligen Lösungen von $\sum k_i v_i = 0$.
Dies ist äquivalent zu einer Frage über die chromatische Zahl eines gewissen
Cayley-Graphen \cite{TothCayley2001}.
\end{bemerkung}

%% ============================================================
\section{Obere Schranken für den Einsamkeitszeitpunkt}\label{sec:obere_schranken}
%% ============================================================

\begin{satz}[Schranke für Einsamkeitszeitpunkt]\label{satz:einsamkeitszeit}
Gilt die LRC für $n$ Läufer mit Geschwindigkeiten $v_1, \ldots, v_n$, so ist
der erste "Einsamkeitszeitpunkt" $t^* > 0$ durch
\[
  t^* \leq \frac{n!}{v_1 \cdots v_n}
\]
beschränkt.
\end{satz}

\begin{proof}[Beweisskizze]
Die "schlechte" Menge für Läufer $i$ besteht aus periodisch wiederkehrenden
Intervallen der Länge $2/(v_i(n+1))$. Das erste $t > 0$, das alle schlechten Mengen
vermeidet, lässt sich durch die $n!$ möglichen Anordnungen der Positionen der
Läufer abschätzen.
\end{proof}

%% ============================================================
\section{Probabilistische Argumente}\label{sec:probabilistisch}
%% ============================================================

\begin{satz}[Zufällige Geschwindigkeiten]\label{satz:random}
Werden $v_1, \ldots, v_n$ unabhängig und gleichmäßig aus $[1, M]$ gewählt, so
geht die Wahrscheinlichkeit, dass die LRC scheitert (kein Einsamkeitszeitpunkt
existiert), für $M \to \infty$ gegen 0.
\end{satz}

\begin{proof}[Beweisskizze]
Der Erwartungswert des Maßes der "stets blockierten" Menge ist durch
$\left(\frac{2}{n+1}\right)^n$ beschränkt. Durch die Markov-Ungleichung ist die
Wahrscheinlichkeit einer nichtleeren "stets blockierten" Menge klein.
\end{proof}

%% ============================================================
\section{Neuere Fortschritte}\label{sec:fortschritte}
%% ============================================================

\begin{satz}[Czerwiński \cite{Czerwinski2012}]\label{satz:czerwinski}
Für $n$ Läufer, deren Geschwindigkeiten eine geometrische Progression
(Verhältnis $q \geq 2$) bilden, gilt die LRC.
\end{satz}

\begin{satz}[Perarnau--Serra \cite{PerarnauSerra2016}]\label{satz:perarnau}
Gilt $\max v_i / \min v_i \leq n+1$, so gilt die LRC.
\end{satz}

\begin{bemerkung}
Diese Ergebnisse decken wichtige Spezialfälle ab, lassen die allgemeine Vermutung
aber offen. Keine bekannte Beweistechnik reicht bislang über $n = 7$ hinaus.
\end{bemerkung}

%% ============================================================
\section{Offene Probleme}\label{sec:offen}
%% ============================================================

\begin{vermutung}[Lonely-Runner-Vermutung -- Allgemeiner Fall]\label{verm:lrc_allg}
Für alle $n \geq 1$ und alle paarweise verschiedenen positiven ganzen Zahlen
$v_1, \ldots, v_n$ gibt es $t \in \R$ mit $\|v_i t\| \geq 1/(n+1)$ für alle $i$.
Dies ist die Hauptvermutung; sie bleibt für $n \geq 8$ offen.
\end{vermutung}

\begin{vermutung}[Starke LRC]\label{verm:starke_lrc}
Für alle $n$ und alle Geschwindigkeiten $v_1, \ldots, v_n$ hat die "gute" Menge
$\{t : \|v_i t\| \geq 1/(n+1) \text{ für alle } i\}$ positives Maß.
\end{vermutung}

\begin{vermutung}[Quantitative LRC]\label{verm:quant_lrc}
Das Maß der guten Menge ist mindestens $(1/(n+1))^n$, mit Gleichheit für die
arithmetischen Geschwindigkeiten $v_i = i$.
\end{vermutung}

\begin{vermutung}[Einsamkeitszahl]\label{verm:einsamkeitszahl}
Definiere die \emph{Einsamkeitszahl} $\lambda(n)$ als das Infimum aller $\alpha$,
sodass für alle $n$ Geschwindigkeiten ein $t$ mit $\|v_i t\| \geq \alpha$ für alle $i$
existiert. Die LRC behauptet $\lambda(n) \geq 1/(n+1)$.
Ob $\lambda(n) = 1/(n+1)$ oder $\lambda(n) > 1/(n+1)$ für großes $n$, ist offen.
\end{vermutung}

\begin{bemerkung}
Der Fall $n = 8$ wurde von mehreren Gruppen mit computerunterstützten Methoden
behauptet, aber kein vollständig verifizierter Beweis wurde bis 2026 akzeptiert.
\end{bemerkung}

%% ============================================================
\section{Schluss}
%% ============================================================

Die Lonely-Runner-Vermutung ist ein täuschend einfaches Problem der kombinatorischen
Zahlentheorie und der diophantischen Approximation. Ihre Äquivalenz zum
Sichtverstopfungsproblem und zu chromatischen Zahlfragen für Cayley-Graphen offenbart
tiefe Verbindungen zu verschiedenen Gebieten der Mathematik. Trotz Beweisen für
$n \leq 7$ Läufer und viele spezielle Geschwindigkeitskonfigurationen (arithmetische
Progressionen, geometrische Progressionen, beschränkte Verhältnisse) bleibt die
allgemeine Vermutung weit offen.

%% Literaturverzeichnis
\begin{thebibliography}{99}

\bibitem{BarajasSerra2008}
J.~Barajas, O.~Serra,
\textit{The lonely runner with seven runners},
Electron.\ J.\ Combin.\ \textbf{15} (2008), R48.

\bibitem{BetkeWills1972}
U.~Betke, J.\,M.~Wills,
\textit{Untere Schranken für zwei diophantische Approximations-Funktionen},
Monatsh.\ Math.\ \textbf{76} (1972), 214--217.

\bibitem{BieniaEtal1998}
W.~Bienia, L.~Goddyn, P.~Gvozdjak, A.~Gyárfás, M.~Tarsi,
\textit{Flows, view obstructions, and the lonely runner},
J.\ Combin.\ Theory Ser.\ B \textbf{72} (1998), 1--9.

\bibitem{BohmanHolzmanKleitman2001}
T.~Bohman, R.~Holzman, D.~Kleitman,
\textit{Six lonely runners},
Electron.\ J.\ Combin.\ \textbf{8} (2001), R3.

\bibitem{Cusick1973}
T.\,W.~Cusick,
\textit{View-obstruction problems},
Aequationes Math.\ \textbf{9} (1973), 165--170.

\bibitem{CusickPomerance1984}
T.\,W.~Cusick, C.~Pomerance,
\textit{View-obstruction problems. III},
J.\ Number Theory \textbf{19} (1984), 131--139.

\bibitem{Czerwinski2012}
S.~Czerwi\'nski,
\textit{Lonely runners for integer linear progressions},
Preprint, 2012.

\bibitem{PerarnauSerra2016}
G.~Perarnau, O.~Serra,
\textit{Correlation among runners and some results on the Lonely Runner Conjecture},
Electron.\ J.\ Combin.\ \textbf{23} (2016), P1.50.

\bibitem{TothCayley2001}
C.\,D.~Tóth,
\textit{The Lonely Runner Conjecture and the Cayley graph},
Manuskript, 2001.

\bibitem{Wills1967}
J.\,M.~Wills,
\textit{Zwei Sätze über inhomogene diophantische Approximation von Irrationalzahlen},
Monatsh.\ Math.\ \textbf{71} (1967), 263--269.

\end{thebibliography}

\end{document}
